Un film de savon d'indice de réfraction \num{\v{n}} et d'épaisseur \qty{\v{ep}}{nm} est entouré d'air. On l'éclaire en lumière blanche et on l'observe en réflexion, sous incidence normale.

\begin{enumerate}[label=\alph*)]
    \item Déterminez les trois plus grandes longueurs d'onde (dans l'air) qui sont fortement réfléchies par le film. Lesquelles font partie du visible (\qtyrange{400}{700}{nm}) ?
    \item Quelle est la plus grande longueur d'onde complètement absente de la lumière réfléchie ?
    \item Quelle épaisseur minimale (non nulle) devrait avoir un film de ce savon pour réfléchir fortement la lumière de \qty{\v{lam}}{nm} ?
\end{enumerate}
